The arithmetic-geometric mean provides a swift way to calculate elliptic integrals numerically.
Its formulation seems rather arbitrary, but we will quickly see that there is a foundational connection between them.
Further reading is highly encouraged, in particular \emph{Pi and the AGM}\footnote{\cite{borwein1987pi} \textsc{Borwein} \& \textsc{Borwein}} by the \textsc{Borwein} brothers.
Following the literature, we define the arithmetic geometric mean $M$ of $a$ and $b$ through the iterative process
\[
a_{n} = \frac{a_{n-1} + b_{n-1}}{2}, \qquad b_{n} = \sqrt{a_{n-1} b_{n-1}}.
\]
This process is maintained until $c_{n} = \sfrac{1}{2}(a_{n-1} - b_{n-1}) < \varepsilon$, where $\varepsilon$ is the desired accuracy.
We have assumed that $0 < b_0 \leq a_0$, and we may observe that it is always true that $b_{n-1} \leq b_n \leq a_n \leq a_{n-1}$, so that the values have a common limit.
This limit is the arithmetic-geometric mean:
\[
M(a, b) = \lim_{n \to \infty} a_n = \lim_{n \to \infty} b_n.
\]
It is evident that from this definition, the starting the algorithm at some later iteration must result in the same function.
In other words,
\[
M(a, b) = M{\left( \frac{a + b}{2}, \sqrt{ab} \right)}.
\]
The arithmetic-geometric mean is also obviously homogeneous, meaning $M(\lambda a, \lambda b) = \lambda M(a,b)$, where $\lambda > 0$.

The choice of $a$, $b$, and $c$ is not necessarily coincidental.
Normalizing the ellipse with respect to the major semiaxis, we get minor semiaxis $\sfrac{b}{a}$, and focus $\sqrt{1 - \sfrac{b^2}{a^2}}$.
It may be shown that\footnote{\cite{borwein1987pi} \textsc{Borwein} \& \textsc{Borwein}, pp.5--7}
\[
\frac{\pi}{2 M(a, b)} = \int_{0}^{\sfrac{\pi}{2}} \frac{1}{\sqrt{a^2 \cos^2{\theta} + b^2 \sin^2{\theta}}} \,\dee\theta.
\]
Factoring out $a$, and using the Pythagorean identity $\cos^2{\theta} = 1 - \sin^2{\theta}$, we may insert the normalized focus, allowing us to identify
\[
k\prime = \sfrac{b}{a}, \qquad k = \sqrt{1 - \sfrac{b^2}{a^2}}
\]
by homogeneity.
The relation to elliptic integrals is now quite apparent, in particular
\[
\ellK(k) = \frac{\pi}{2 M(1, k^{\prime})}, \quad \frac{\ellE(k)}{\ellK(k)} = 1 - \sum_{n} 2^{n-1} {c_n}^2.
\]
These are the \texttt{K} and \texttt{E} methods implemented in the \texttt{Jacobi} class of \texttt{jacobi.py}.
Writing now that ${c_0}^2 = {a_0}^2 - {b_0}^2$, we want to find the natural choices for $a_0$ and $b_0$, for calculating $\ellK(k)$.
It is clear that ${c_n}^2 = {a_n}^2 - {b_n}^2$, and in turn that $1 = \sfrac{{b_n}^2}{{a_n}^2} + \sfrac{{c_n}^2}{{a_n}^2}$.
This we identify with $1 = {k\prime}^2 + k^2$, yielding $a_0 = 1$ and $b_0 = k\prime$ for calculating the elliptic integral.

At the very heart of the implementation of the arithmetic-geometric mean is the descending \textsc{Landen} transform.\footnote{\cite{king1924direct} \textsc{King}, sect.II, IV}
The premise is to start at some modulus $k$ and amplitude $\varphi$, and set
\begin{equation}\label{eq:descending_landen_transform}
k_1 = \frac{1 - k^{\prime}}{1 + k^{\prime}}, \quad \sin{\varphi_1} = \frac{(1 + k^{\prime}) \sin{\varphi} \cos{\varphi}}{\sqrt{1 - k^2 \sin^2{\varphi}}},
\end{equation}
where $\varphi \leq \varphi_1$.
It can then be shown that
\[
(1 + k^{\prime}) \ellF(\varphi \vert k) = \ellF(\varphi_1 \vert k_1).
\]
Furthermore, it can be shown that equation \eqref{eq:descending_landen_transform} is equivalent to stating that
\[
\sin{( 2\varphi - \varphi_1 )} = k_1 \sin{\varphi_1}.
\]
Through some tedious algebra, we find that this can be written now in terms of an arithmetic-geometric mean algorithm,
\begin{equation}\label{eq:AGM_phi}
\varphi_n = \sfrac{1}{2} \arcsin{\left( \frac{c_{n+1}}{a_{n+1}} \sin{\varphi_{n+1}} \right)} + \frac{\varphi_{n+1}}{2}.
\end{equation}
The above implies that if we run the arithmetic-geometric mean algorithm up until $c_N$ is sufficiently small, we can iterate backwards to find a $\varphi_0$, whence the \emph{descending} \textsc{Landen} transform.
We obviously know what parameters to choose for the arithmetic-geometric mean algorithm, but the starting value of $\varphi_N$ is not as obvious.
It can be shown that the following holds.\footnote{\cite{king1924direct} \textsc{King}, eq.25}
\begin{equation}\label{eq:AGM_ellF}
\ellF(\varphi \vert k) = \frac{\varphi_N}{a_N 2^N}
\end{equation}
This comes in very handy should we want to calculate any of the \textsc{Jacobi} elliptic functions, since $\ellF$ is the argument.
The angle may then be found with the algorithm in equation \eqref{eq:AGM_phi}, which is what the $\phitt$ method of the \texttt{Jacobi} class does.
